\documentclass[12pt,a4paper,oneside]{article}
\usepackage[brazil]{babel}
\usepackage[utf8]{inputenc}
\usepackage[dvipsnames]{xcolor}
\usepackage{listings}
\usepackage{wrapfig}
\usepackage{olymp}

\setcounter{problem}{1}

\begin{document}

\begin{problem}{Cotação do dólar}{dolar}{2}

\begingroup
\setlength{\intextsep}{0pt}%
% \setlength{\columnsep}{0pt}
\begin{wrapfigure}{r}{0.35\textwidth}
\includegraphics[width=0.95\linewidth]{images/exercise_102_0.png} 
%\caption{Caption1}
%\label{fig:subim1}
\end{wrapfigure}

Na semana passada o dólar comercial fechou cotado para compra em mais de R\$5,70 na compra, o maior valor de outubro. Na semana anterior chegou a ser cotado abaixo de R\$5,50. O gráfico ao lado ilustra a variação nos últimos dias.

Embora uma análise dessa situação seja importante sob vários aspectos econômicos, nesta questão estamos interessados apenas no valor de alta e de queda do dólar em algum período. Particularmente, dadas as cotações de vários dias de semana consecutivos, queremos saber qual a maior alta e a maior queda neste período.

Por exemplo, considere a cotação nos últimos 7 dias do gráfico:

{%\small
\texttt{
5.5565
5.6417
5.7111
5.5967
5.5794
5.5661
5.6118
}
}

A maior queda foi de \texttt{0.1144} (de \texttt{5.7111} para \texttt{5.5967}).

A maior alta foi de \texttt{0.0852} (de \texttt{5.5565} para \texttt{5.6417}).

% \Notes

% Uma seleção ganha 3 pontos para cada vitória e 1 ponto para cada empate. Derrotas não alteram sua pontuação.

\InputFile

A entrada começa com uma linha contendo um número inteiro $N$, a quantidade de cotações. Em seguida $N$ linhas, cada uma contendo um número real, indicando a cotação de cada dia.

Restrições: $4 \leq N \leq 100$, cada cotação entre $1.0000$ e $6.0000$.
\endgroup

\OutputFile

Imprima duas linhas na saída, a primeira com o valor da maior queda e a segunda com o valor da maior alta. Os dois valores devem conter 4 casas decimais. Sempre haverá pelo menos uma queda e pelo menos uma alta nas cotações da entrada.

% \Notes
% Preste atenção aos tipos das variáveis, pois $F(90) \approx 2^{61}$.

\Example

\begin{example}
\exmp{
4
2.0000
3.0000
4.0000
2.0000
}{
2.0000
1.0000
}%
\end{example}

\begin{example}
\exmp{
7
5.5565
5.6417
5.7111
5.5967
5.5794
5.5661
5.6118
}{
0.1144
0.0852
}%
\end{example}


%Comentários sobre o exemplo, se necessário.

\end{problem}

\end{document}
